% Options for packages loaded elsewhere
\PassOptionsToPackage{unicode}{hyperref}
\PassOptionsToPackage{hyphens}{url}
\PassOptionsToPackage{dvipsnames,svgnames,x11names}{xcolor}
%
\documentclass[
  11pt,
]{article}
\usepackage{amsmath,amssymb}
\usepackage{iftex}
\ifPDFTeX
  \usepackage[T1]{fontenc}
  \usepackage[utf8]{inputenc}
  \usepackage{textcomp} % provide euro and other symbols
\else % if luatex or xetex
  \usepackage{unicode-math} % this also loads fontspec
  \defaultfontfeatures{Scale=MatchLowercase}
  \defaultfontfeatures[\rmfamily]{Ligatures=TeX,Scale=1}
\fi
\usepackage[]{charter}
\ifPDFTeX\else
  % xetex/luatex font selection
\fi
% Use upquote if available, for straight quotes in verbatim environments
\IfFileExists{upquote.sty}{\usepackage{upquote}}{}
\IfFileExists{microtype.sty}{% use microtype if available
  \usepackage[]{microtype}
  \UseMicrotypeSet[protrusion]{basicmath} % disable protrusion for tt fonts
}{}
\makeatletter
\@ifundefined{KOMAClassName}{% if non-KOMA class
  \IfFileExists{parskip.sty}{%
    \usepackage{parskip}
  }{% else
    \setlength{\parindent}{0pt}
    \setlength{\parskip}{6pt plus 2pt minus 1pt}}
}{% if KOMA class
  \KOMAoptions{parskip=half}}
\makeatother
\usepackage{xcolor}
\usepackage[margin=2.5cm]{geometry}
\setlength{\emergencystretch}{3em} % prevent overfull lines
\providecommand{\tightlist}{%
  \setlength{\itemsep}{0pt}\setlength{\parskip}{0pt}}
\setcounter{secnumdepth}{-\maxdimen} % remove section numbering
% ==========================================================================
%  IA na Sociedade — cabeçalho LaTeX (pdflatex)
%  Programa CiberExt 26-29 · FEELT38103
%  Adaptado de "AI in Society" (Universidade de Helsinque / Una Europa)
%  Licença: CC BY-NC-SA 4.0
% --------------------------------------------------------------------------
%  Espelha a identidade do ia-style.css: papel quente, tinta ardósia e a
%  faixa tricolor dos três eixos do curso (conceitos / aplicações / riscos).
% ==========================================================================

% ---------- Fontes ----------
% A serifada vem do build.sh, que passa  -V fontfamily=charter  ao pandoc.
% Aqui só acrescentamos a sem-serifa; ambas com fallback, porque nem toda
% instalação de LaTeX traz os mesmos pacotes.
\IfFileExists{helvet.sty}{\usepackage[scaled=0.90]{helvet}}{}
\usepackage{textcomp}
\IfFileExists{microtype.sty}{\usepackage{microtype}}{}

% ---------- Pacotes de layout ----------
\usepackage{xcolor}
\usepackage{tcolorbox}
\tcbuselibrary{skins,breakable}
\usepackage{titlesec}
\usepackage{fancyhdr}
\usepackage{enumitem}
\usepackage{amssymb}
\usepackage{booktabs}
\usepackage{array}
\usepackage{colortbl}
\usepackage{ragged2e}
\usepackage{fvextra}
\usepackage{newunicodechar}

% ---------- Caracteres Unicode que o pdflatex não conhece ----------
\newunicodechar{✓}{{\color{iaAplic}\checkmark}}
\newunicodechar{✅}{{\color{iaAplic}\checkmark}}
\newunicodechar{→}{\ensuremath{\rightarrow}}
\newunicodechar{←}{\ensuremath{\leftarrow}}
\newunicodechar{↔}{\ensuremath{\leftrightarrow}}
\newunicodechar{–}{--}
\newunicodechar{—}{---}
\newunicodechar{€}{\texteuro}
\newunicodechar{£}{\textsterling}
\newunicodechar{×}{\ensuremath{\times}}
\newunicodechar{≈}{\ensuremath{\approx}}
\newunicodechar{≥}{\ensuremath{\geq}}
\newunicodechar{≤}{\ensuremath{\leq}}
\newunicodechar{•}{\textbullet}
\newunicodechar{″}{''}
\newunicodechar{′}{'}

% ==========================================================================
%  PALETA — os mesmos valores do ia-style.css
% ==========================================================================
\definecolor{iaConceitos}{HTML}{24505C}   % petróleo — títulos, eixo 1
\definecolor{iaAplic}{HTML}{16746C}       % verdete  — subtítulos, eixo 2
\definecolor{iaRiscos}{HTML}{B45F14}      % açafrão  — destaques, eixo 3
\definecolor{iaAmeixa}{HTML}{77345F}      % links internos, termos técnicos
\definecolor{iaTinta}{HTML}{1F2933}       % corpo do texto
\definecolor{iaTintaFraca}{HTML}{5C6670}
\definecolor{iaRegra}{HTML}{E4DBCD}
\definecolor{iaBloco}{HTML}{F3EEE4}
\definecolor{iaNota}{HTML}{EAF1F0}
\definecolor{iaCaso}{HTML}{FCF3E6}
\definecolor{iaFilos}{HTML}{F7EFF4}

\color{iaTinta}

% ==========================================================================
%  ASSINATURA — faixa tricolor dos três eixos
% ==========================================================================
\newcommand{\faixaEixos}[1][\linewidth]{%
  \noindent
  \textcolor{iaConceitos}{\rule{0.3333#1}{4pt}}%
  \textcolor{iaAplic}{\rule{0.3333#1}{4pt}}%
  \textcolor{iaRiscos}{\rule{0.3333#1}{4pt}}%
  \par}

% ---------- Página de título ----------
\makeatletter
\newcommand{\subtitle}[1]{\gdef\@subtitle{#1}}
\providecommand{\@subtitle}{}
\def\@maketitle{%
  \newpage\null\vskip 1.5em%
  \faixaEixos
  \vskip 2.2em
  {\raggedright
    {\sffamily\bfseries\huge\color{iaConceitos}\@title\par}%
    \vskip 0.7em%
    {\sffamily\large\color{iaTintaFraca}\@subtitle\par}%
    \vskip 2em%
    {\color{iaRegra}\rule{\linewidth}{1pt}}%
    \vskip 1em%
    {\sffamily\small\color{iaTinta}\@author\par}%
    \vskip 0.5em%
    {\sffamily\footnotesize\color{iaTintaFraca}\@date\par}%
  }\par\vskip 2.5em}
\makeatother

% ==========================================================================
%  TÍTULOS DE SEÇÃO
% ==========================================================================
\titleformat{\section}[block]
  {\sffamily\Large\bfseries\color{iaConceitos}}
  {\thesection}{0.7em}{}
  [{\vspace{2pt}\color{iaRegra}\titlerule[1.2pt]}]
\titleformat{\subsection}
  {\sffamily\large\bfseries\color{iaAplic}}
  {\thesubsection}{0.6em}{}
\titleformat{\subsubsection}
  {\sffamily\normalsize\bfseries\color{iaRiscos}}
  {\thesubsubsection}{0.5em}{}
\titleformat{\paragraph}[runin]
  {\sffamily\small\bfseries\color{iaRiscos}}
  {}{0em}{}

\titlespacing*{\section}{0pt}{2.6ex plus 1ex minus .2ex}{1.3ex plus .2ex}
\titlespacing*{\subsection}{0pt}{2.2ex plus 1ex minus .2ex}{1ex plus .2ex}

% ==========================================================================
%  CABEÇALHO E RODAPÉ
% ==========================================================================
\setlength{\headheight}{24pt}
\renewcommand{\sectionmark}[1]{\markboth{#1}{}}
\pagestyle{fancy}
\fancyhf{}
\renewcommand{\headrulewidth}{0.6pt}
\renewcommand{\footrulewidth}{0pt}
\fancyhead[L]{\sffamily\footnotesize\color{iaAplic}Ética da IA}
\fancyhead[R]{\sffamily\footnotesize\color{iaAplic}\nouppercase{\leftmark}}
\fancyfoot[C]{\sffamily\footnotesize\color{iaConceitos}\thepage}
\renewcommand{\headrule}{%
  \hbox to\headwidth{\color{iaRegra}\leaders\hrule height \headrulewidth\hfill}}

% ==========================================================================
%  TEXTO CORRIDO
% ==========================================================================
\linespread{1.08}
\setlength{\parskip}{0.55em}
\setlength{\parindent}{0pt}
\emergencystretch=3em
\hbadness=10000

% ---------- Termos técnicos / código inline ----------
\let\oldtexttt\texttt
\renewcommand{\texttt}[1]{{\color{iaAmeixa}\small\ttfamily #1}}

% ---------- Blocos verbatim ----------
\fvset{breaklines=true,breakanywhere=true,fontsize=\small}

% ---------- Blocos de código (ambiente Shaded do pandoc) ----------
\renewenvironment{Shaded}
  {\begin{tcolorbox}[
      breakable, enhanced jigsaw, rounded corners,
      arc=3pt, boxrule=0pt, frame hidden,
      colback=iaBloco,
      borderline west={3pt}{0pt}{iaAmeixa},
      left=10pt, right=8pt, top=6pt, bottom=6pt,
      before skip=9pt, after skip=9pt]}
  {\end{tcolorbox}}

% ---------- Citações: caixa de caso ----------
\renewenvironment{quote}
  {\begin{tcolorbox}[
      breakable, enhanced jigsaw, rounded corners,
      arc=3pt, boxrule=0pt, frame hidden,
      colback=iaCaso,
      borderline west={3pt}{0pt}{iaRiscos},
      left=12pt, right=10pt, top=9pt, bottom=9pt,
      before skip=9pt, after skip=9pt]}
  {\end{tcolorbox}}

% ==========================================================================
%  BLOCOS DE APOIO
%  No markdown, escreva:   ::: nota   ...texto...   :::
%  O pandoc converte fenced_divs para os ambientes abaixo.
% ==========================================================================
\newtcolorbox{iabox}[3]{
  breakable, enhanced jigsaw, rounded corners,
  arc=3pt, boxrule=0pt, frame hidden,
  colback=#2,
  borderline west={3pt}{0pt}{#1},
  left=12pt, right=10pt, top=9pt, bottom=9pt,
  before skip=11pt, after skip=11pt,
  before upper={\setlength{\parskip}{0.5em}\setlength{\parindent}{0pt}},
  title={\sffamily\bfseries\footnotesize\textcolor{#1}{\MakeUppercase{#3}}},
  attach title to upper={\par\vspace{5pt}},
  fonttitle=\normalfont}

% Cada ambiente aceita um titulo opcional entre colchetes, vindo do
% atributo data-titulo do markdown:   \begin{filosofia}[Titulo] ... \end{filosofia}
% Os quatro primeiros espelham os icones do material original.
\newenvironment{filosofia}[1][Conceito filosofico]{\begin{iabox}{iaAmeixa}{iaFilos}{#1}}{\end{iabox}}
\newenvironment{tecnica}[1][Nota tecnica]{\begin{iabox}{iaConceitos}{iaNota}{#1}}{\end{iabox}}
\newenvironment{contexto}[1][Contexto]{\begin{iabox}{iaAplic}{iaBloco}{#1}}{\end{iabox}}
\newenvironment{definicao}[1][Definicao]{\begin{iabox}{iaRiscos}{iaCaso}{#1}}{\end{iabox}}

\newenvironment{nota}[1][Nota]{\begin{iabox}{iaConceitos}{iaNota}{#1}}{\end{iabox}}
\newenvironment{caso}[1][Caso]{\begin{iabox}{iaRiscos}{iaCaso}{#1}}{\end{iabox}}
\newenvironment{reflexao}[1][Para pensar]{\begin{iabox}{iaAplic}{iaBloco}{#1}}{\end{iabox}}
\newenvironment{licenca}
  {\par\vspace{2em}\faixaEixos\par\vspace{0.8em}
   \begingroup\sffamily\scriptsize\color{iaTintaFraca}}
  {\endgroup\par}

% ==========================================================================
%  LISTAS
% ==========================================================================
\setlist[itemize]{leftmargin=1.4em, itemsep=2pt, topsep=4pt, parsep=0pt}
\setlist[enumerate]{leftmargin=1.7em, itemsep=2pt, topsep=4pt, parsep=0pt}
\renewcommand{\labelitemi}{\textcolor{iaRiscos}{\textbullet}}
\renewcommand{\labelitemii}{\textcolor{iaAplic}{--}}

% ==========================================================================
%  TABELAS
% ==========================================================================
\renewcommand{\arraystretch}{1.3}
\arrayrulecolor{iaRegra}

% ==========================================================================
%  RÓTULOS EM PORTUGUÊS
%  Definidos aqui (e não via -V na linha de comando) porque acentos passados
%  como argumento dependem do locale do terminal e podem se corromper.
% ==========================================================================
\renewcommand{\contentsname}{Sumário}
\renewcommand{\refname}{Referências}
\renewcommand{\figurename}{Figura}
\renewcommand{\tablename}{Tabela}
\ifLuaTeX
  \usepackage{selnolig}  % disable illegal ligatures
\fi
\IfFileExists{bookmark.sty}{\usepackage{bookmark}}{\usepackage{hyperref}}
\IfFileExists{xurl.sty}{\usepackage{xurl}}{} % add URL line breaks if available
\urlstyle{same}
\hypersetup{
  pdftitle={Ética da Inteligência Artificial},
  colorlinks=true,
  linkcolor={iaAmeixa},
  filecolor={Maroon},
  citecolor={Blue},
  urlcolor={iaRiscos},
  pdfcreator={LaTeX via pandoc}}

\title{Ética da Inteligência Artificial}
\usepackage{etoolbox}
\makeatletter
\providecommand{\subtitle}[1]{% add subtitle to \maketitle
  \apptocmd{\@title}{\par {\large #1 \par}}{}{}
}
\makeatother
\subtitle{Capítulo 2 --- Exercícios}
\author{Tradução e adaptação para o português: José Lucas Lira Bizil,
Fernando Mazzeto Lisboa Lima e Matheus da Silva Fernandes\\
Programa CiberExt 26-29 · FEELT38103 · Universidade Federal de
Uberlândia\\
Original: \emph{Ethics of AI}, Universidade de Helsinque}
\date{2026}

\begin{document}
\maketitle

{
\hypersetup{linkcolor=iaConceitos}
\setcounter{tocdepth}{2}
\tableofcontents
}
\hypertarget{capuxedtulo-2-exercuxedcios}{%
\section{Capítulo 2 --- Exercícios}\label{capuxedtulo-2-exercuxedcios}}

\hypertarget{exercuxedcio-4-dos-benefuxedcios-imediatos-aos-objetivos-uxe9ticos-de-longo-prazo-dissertativa}{%
\subsection{\texorpdfstring{Exercício 4 --- Dos benefícios imediatos aos
objetivos éticos de longo prazo
\texttt{{[}dissertativa{]}}}{Exercício 4 --- Dos benefícios imediatos aos objetivos éticos de longo prazo {[}dissertativa{]}}}\label{exercuxedcio-4-dos-benefuxedcios-imediatos-aos-objetivos-uxe9ticos-de-longo-prazo-dissertativa}}

\textbf{Pontuação:} 3 pontos · \textbf{Extensão:} 75 a 1000 palavras ·
\textbf{Tentativas:} ilimitadas \textbf{Componente Open edX:}
\texttt{openassessment} (revisão por pares)

Lembre-se do exemplo que abriu este capítulo, sobre a migração de uma
saúde reativa para uma saúde preventiva. Como diretor digital, pede-se
agora que você tome uma decisão. A organização de saúde da cidade
deveria migrar da saúde reativa para a preventiva e começar a implantar
métodos de análise de dados?

Os sistemas preveem os possíveis riscos de saúde dos cidadãos analisando
um grande número de critérios a partir de diversos prontuários médicos.
São projetados para identificar indivíduos de alto risco. Se essas
pessoas puderem ser priorizadas, isso pode permitir melhor estimativa de
impacto e melhor planejamento dos serviços básicos de saúde.

A saúde preventiva também tem potencial de reduzir significativamente os
custos sociais e de saúde. Essas economias, enfatiza o relatório,
poderiam ser usadas para outras finalidades, como o desenvolvimento de
escolas.

Responda brevemente às questões a seguir. Avalie:

\begin{enumerate}
\def\labelenumi{\arabic{enumi}.}
\tightlist
\item
  Os benefícios e riscos \textbf{imediatos}, do aqui e agora, de
  implantar os sistemas de saúde preventiva.
\item
  Os impactos de \textbf{médio prazo} --- como esses sistemas poderiam
  afetar o mercado de trabalho, a educação ou, por exemplo, a
  distribuição do bem-estar entre os cidadãos?
\item
  Os objetivos éticos fundamentais de \textbf{longo prazo} do
  desenvolvimento e uso de análise de dados em saúde preventiva. Como o
  cidadão é visto enquanto cidadão? Quais são os princípios ou valores
  éticos promovidos pela cidade? Como o uso da tecnologia é justificado
  eticamente?
\end{enumerate}

Responda a todos os itens acima. Sua resposta será revisada por outros
participantes e pelos instrutores. Confira antes de enviar: uma vez
enviada, não poderá mais editá-la.

\begin{center}\rule{0.5\linewidth}{0.5pt}\end{center}

\hypertarget{exercuxedcio-5-como-a-ia-poderia-promover-o-bem-dissertativa-quatro-partes}{%
\subsection{\texorpdfstring{Exercício 5 --- Como a IA poderia promover o
bem?
\texttt{{[}dissertativa\ ·\ quatro\ partes{]}}}{Exercício 5 --- Como a IA poderia promover o bem? {[}dissertativa · quatro partes{]}}}\label{exercuxedcio-5-como-a-ia-poderia-promover-o-bem-dissertativa-quatro-partes}}

\textbf{Pontuação:} 2 pontos · \textbf{Extensão:} 75 a 1000 palavras por
item · \textbf{Tentativas:} ilimitadas \textbf{Componente Open edX:}
\texttt{openassessment} com quatro campos de resposta

Segundo o AI4People, a IA poderia oferecer ``quatro grandes
oportunidades para a sociedade'' se for desenvolvida de modo a respeitar
a dignidade humana. Nesta tarefa, seu trabalho é descobrir como essas
oportunidades poderiam ser alcançadas na prática.

\hypertarget{autorrealizauxe7uxe3o}{%
\subsubsection{1. Autorrealização}\label{autorrealizauxe7uxe3o}}

Significa a capacidade das pessoas de florescer em termos de seus
próprios projetos de vida, interesses, potencialidades e habilidades.
Segundo o AI4People, isso pode ser alcançado se a automação
``inteligente'' for usada para libertar as pessoas do trabalho
rotineiro. Contudo, o risco nesse caso é o ritmo do desenvolvimento
tecnológico e a possível distribuição desigual dos custos e benefícios
resultantes. Uma desvalorização muito rápida pode, por exemplo, levar a
uma ruptura acelerada do mercado de trabalho, mudando a natureza do
emprego e causando sério dano econômico.

\begin{quote}
\textbf{Na prática, como a automação inteligente poderia ser usada de
modo a libertar as pessoas do trabalho rotineiro? Forneça dois exemplos
concretos, com breve explicação de como, exatamente, a automação liberta
as pessoas.}
\end{quote}

\hypertarget{ampliauxe7uxe3o-da-aguxeancia-humana}{%
\subsubsection{2. Ampliação da agência
humana}\label{ampliauxe7uxe3o-da-aguxeancia-humana}}

A IA pode fornecer um reservatório crescente de ``agência inteligente'',
o que significa que pode ampliar a inteligência humana e nos dar enorme
apoio cognitivo. Isso, porém, exige uma compreensão clara de
responsabilidade: que tipo de IA desenvolvemos, como a usamos e se
compartilhamos suas vantagens e benefícios com todos.

\begin{quote}
\textbf{De que maneiras a IA pode ampliar a inteligência humana ou
apoiar nosso pensamento? Forneça dois exemplos e explique como a IA nos
ajuda a pensar.}
\end{quote}

\hypertarget{desenvolvimento-de-capacidades}{%
\subsubsection{3. Desenvolvimento de
capacidades}\label{desenvolvimento-de-capacidades}}

A IA oferece oportunidades para desenvolver as capacidades dos
indivíduos e da sociedade em geral. O uso de tecnologias de IA apresenta
inúmeras possibilidades de ``reinventar'' a sociedade, ampliando
radicalmente aquilo de que os humanos são coletivamente capazes. Mais IA
pode apoiar melhor coordenação e, portanto, objetivos mais ambiciosos.
Contudo, se recorrermos ao uso de tecnologias de IA para ampliar nossas
próprias capacidades da maneira errada, podemos delegar tarefas
importantes --- e, sobretudo, decisões --- a sistemas autônomos que
deveriam permanecer ao menos parcialmente sujeitos à supervisão e à
escolha humanas.

\begin{quote}
\textbf{Como a IA poderia melhorar a coordenação coletiva das sociedades
no futuro? Forneça um exemplo e explique brevemente.}
\end{quote}

\hypertarget{cultivo-da-coesuxe3o-social}{%
\subsubsection{4. Cultivo da coesão
social}\label{cultivo-da-coesuxe3o-social}}

Muitos problemas globais --- mudança climática, pobreza --- exigem
coordenação complexa. Por exemplo, os esforços para enfrentar a mudança
climática expuseram o desafio de criar uma resposta coesa, tanto dentro
das sociedades quanto entre elas. A IA, com suas soluções intensivas em
dados e movidas por algoritmos, pode ajudar enormemente a lidar com essa
complexidade de coordenação, apoiando mais coesão e colaboração sociais.
Ainda assim, o risco é que o poder computacional e preditivo da IA deve
ser usado para encorajar a autodeterminação humana e fomentar a coesão
social, não para minar a dignidade ou o florescimento humano.

\begin{quote}
\textbf{Como a IA poderia nos ajudar a colaborar melhor e a lidar com
problemas como mudança climática ou pobreza? Forneça dois exemplos
concretos e explique brevemente.}
\end{quote}

\begin{center}\rule{0.5\linewidth}{0.5pt}\end{center}

\hypertarget{avaliauxe7uxe3o-do-capuxedtulo-pesquisa}{%
\subsection{\texorpdfstring{Avaliação do capítulo
\texttt{{[}pesquisa{]}}}{Avaliação do capítulo {[}pesquisa{]}}}\label{avaliauxe7uxe3o-do-capuxedtulo-pesquisa}}

\textbf{Pontuação:} 1 ponto · \textbf{Tentativas:} ilimitadas

\begin{itemize}
\tightlist
\item
  \textbf{Utilidade deste capítulo} --- nota de 1 a 5
\item
  \textbf{Clareza deste capítulo} --- nota de 1 a 5
\item
  \textbf{Dificuldade deste capítulo} --- nota de 1 a 5
\item
  \textbf{Outros comentários} --- até 600 palavras
\end{itemize}

\begin{center}\rule{0.5\linewidth}{0.5pt}\end{center}

\begin{licenca}

\textbf{Material original:} \emph{Ethics of AI}, Universidade de
Helsinque. Professores responsáveis: Anna-Mari Rusanen e Jukka K.
Nurminen. \url{https://ethics-of-ai.mooc.fi/}. \textbf{Esta versão:}
tradução por José Lucas Lira Bizil, Fernando Mazzeto Lisboa Lima e
Matheus da Silva Fernandes, Programa CiberExt 26-29 (FEELT38103), UFU,
2026. Obra derivada. \textbf{Licença:} CC BY-NC-SA 4.0.

\end{licenca}

\end{document}
